\unnumbered{Заключение}
В одной симуляционной постановке исследованы классический подъём со стабилизацией и четыре метода обучения с подкреплением. Явно зафиксированы ускорительный вход, рабочие ограничения, интервал управления, награда и физический критерий.

Основная серия отвечает на вопрос о достижимости рабочего решения при выбранном бюджете. SAC-on и TQC-on дают 20/20 на validation для best и last всех трёх seeds. DQN способен находить успешное решение, но отдельные seeds и последние веса хуже. DDPG-on успешен только на одном исходном seed, SAC-off с базовыми настройками не достигает успеха. Эти результаты характеризуют конкретные процедуры на исходной узкой области; validation использована при выборе и не является независимым итоговым тестом.

Однофакторный подбор показывает ограниченность улучшения настройки. Warmup5000 строго выигрывает Stage 1 для DDPG-on и SAC-off, но минимум успеха по seeds остаётся нулевым. Обе настройки не проходят confirmation на seeds 3–5. У SAC-off предположение об улучшении раннего покрытия верха также не получает поддержки по заранее закреплённому критерию.

Сохранённые DDPG-журналы устанавливают конкретное расхождение между наградой и удержанием: верхнее положение сохраняется, но скорость тележки прерывает требуемый интервал. В одном seed это сочетается с ростом суммарной награды. Такая диагностика объясняет, почему визуальное положение маятника и скалярная награда недостаточны, но не устанавливает причину деградации параметров нейросети.

Завершённая структурированная проверка начальных условий оценивает новые старты без новых обучений. На nominal20 фиксированные SAC/TQC и классика почти полностью воспроизводят решение исходной задачи; на осевых, граничных и совместных отклонениях возникают различия между политиками, скрытые прежней валидацией. На всей фиксированной матрице фильтр повышает успех с 441/980 до 796/980 и устраняет наблюдённые превышения 0,24 м с учётом допуска. При этом 184/980 эпизода on достигают горизонта без этих превышений, но не выполняют заключительное удержание. Измеренный парный эффект относится к 140 выбранным состояниям из $K$ и закреплённым контроллерам, а не ко всей области $K$; это не доказательство универсальной безопасности или превосходства на произвольных состояниях.

Практический результат работы — воспроизводимая картина сильных и слабых сторон конкретных процедур, включая отрицательные исходы. Она позволяет обоснованно отделить задачу защиты тележки от задачи захвата маятника и повторяемость обучения от чувствительности фиксированной политики к старту. Неизвестными остаются общая область захвата, распределение качества по будущим seeds, устойчивость к неопределённости модели и перенос на оборудование. Такие выводы требуют новых данных.

Перспективы работы — независимая оценка зафиксированной процедуры на новых стартах и большем числе обучений, перепланирование классического подъёма, а также учёт задержек, шумов и ошибок исполнения ускорения. Закрытый набор итоговой оценки не использовался. Перенос на оборудование требует отдельной проверки модели и привода.
